\subsection{Smoothing property and existence of limits}


Now we use the lemmas obtained in the previous subsection on our function on curves $f$ satisfying smoothing, convex union, homogeneity and stability to prove the existence of certain limits that will be crucial to construct our measure.

An immediate consequence of smoothing and convex union is positivity of $f$.

\begin{corollary}
  Let $f$ be a functional satisfying convex union and smoothing. Then $f([a]) \geq 0$ for any $a \in \pi_1(S)$.
  \label{cor:positivity}
\end{corollary}
\begin{proof}
Let $x \in \pi_1(S)$ so that $a$ crosses $x$ to the right. By oriented smoothing,
\[
[a][x] \reducesto [ax].
\]
On the other hand, by Equation~\ref{eq:cross-red3} (which uses unoriented smoothing), we have
\[
[ax][a]\reducesto [x]
\]
Since $f$ satisfies smoothing (oriented and unoriented) as well as convex union, we get
\begin{align*}
    f(a) + f(x) &\geq f([a] [x]) \geq f(ax)\\
     f(ax) + f(a) &\geq f([ax] [a]) \geq f(x).\\   
\end{align*}
Adding both equations, we get $f(a) \geq 0$, as we wanted.
\end{proof}

A similar argument with different assumptions also yields positivity.

\begin{corollary}
  Let $f$ be a functional satisfying convex union and quasi-smoothing and stability/homogeneity. Then $f([a]) \geq 0$ for any $a \in \pi_1(S)$.
  \label{cor:positivity2}
\end{corollary}

\begin{proof}
Let $x \in \pi_1(S)$ so that $a$ crosses $x$ to the right. As in the proof of Corollary~\ref{cor:positivity}, we have
\[
[a^n][x] \reducesto [a^nx].
\]
and
\[
[a^nx][a^n]\reducesto [x]
\]
Since $f$ satisfies quasi-smoothing (oriented and unoriented) as well as convex union, we get
\begin{align*}
    f(a^n) + f(x) &\geq f([a^n] [x]) \geq f(a^nx) - R\\
     f(a^nx) + f(a^n) &\geq f([ax] [a^n]) \geq f(x) - R.\\   
\end{align*}
Adding both equations, we get $f(a^n) \geq -2R$. By stability and homogeneity, this yields $n f(a) \geq -2R$, and taking $n$ large enough we see that $f(a) \geq 0$.
\end{proof}

The following result shows a key limit to construct the measure exists and uses the smoothing lemmas.

\begin{corollary}\label{cor:additivity-limit}
Let $f$ satisfy convex union and smoothing.
  Let $x$ be simple. If $a$ crosses $x$ to the right, then
  \[
    \lim_{n \to \infty} f([a x^n]) - f([x^n])
  \]
  exists.
\end{corollary}

\begin{proof}
  By Eq.~\eqref{eq:cross-red2}, and smoothing and additivity (in fact, convex union is enough), we get%
  \footnote{This is where smoothing as opposed to just
    quasi-smoothing is crucial. See Remark~\ref{rmk:nonconvex} for an example.}
  \[
  f([ax^{n-1}]) + f([ax^{n+1}]) \geq 2 f([ax^n]).
  \]
  Hence, $f([a x^n])$ is a convex function
  of $n$, and thus so is $g(n) \coloneqq f([a x^n]) - f([x^n])$.
 By Lemma~\ref{lem:rhs}, we have that
 \[
 [ax^m][x^n] \reducesto [ax^{n+m}]
 \]
 which for $m=0$ yields
 \[
 [a][x^n] \reducesto [ax^{n}].
 \] 
 Thus, by additivity and oriented smoothing of $f$, we get
 \[
 f([a]) + f([x^n]) \geq f([ax^{n}]).
 \]
 Now, because of equation~(\ref{eq:cross-red3}),
 with $m=0$, we get
 \[
 [ax^n][a] \reducesto [x^{n}].
 \]  
 Therefore, by additivity and unoriented smoothing of $f$, we get
  \[
    -f([a]) \le f([a x^n]) - f([x^n]).
  \]
  Thus, $g$ is a convex function that is bounded as $n \to \infty$, and therefore it has a
  well-defined slope at infinity. Thus, the limit in the statement exists.
\end{proof}

\begin{remark}
\label{rmk:nonconvex}
The limit in Corollary~\ref{cor:additivity-limit} need not exist if one only assumes $f$ satisfies quasi-smoothing.
Consider the once-punctured torus with generating set $\{a,b,a^2\}$, where $a$ and $b$ are the standard generators of $\pi_1(S)$.
Let $f$ be the stable word-length with respect to this generating set.
By picking the right orientations, $a$ crosses $b$ right-handedly, thus $ab$ also crosses $b$ right-handedly by Lemma~~\ref{lem:rhs_prod}.
$f$ satisfies quasi-smoothing instead of smoothing (see \cite[Example~4.10]{MGT19:CurvesToCurrents}) for details). 
The expression in the limit in Corollary~\ref{cor:additivity-limit} for the pair $ab$ and $b$ is
\begin{equation*}
f(ab^{n+1})-f(b^n)=\left\{ \begin{array}{cc} 1+ \frac{n+2}{2} - \frac{n}{2}=2& \mbox{ if } n \mbox{ is even } \\ 1+ \frac{n+1}{2} - \frac{n+1}{2}=1 & \mbox{ if } n \mbox{ is odd } \end{array} \right.
\end{equation*}
so the limit when $n$ goes to infinity doesn't exist.
\end{remark}

\begin{proof}[Proof of Proposition~\ref{prop:limitexists}]
This follows immediately from Corollary~\ref{cor:additivity-limit}.
\end{proof}

\begin{remark}
Convexity of $f([ax^n])$ as a function of $n$ also plays a key role in the work of Feng Luo of characterization of intersection numbers for simple closed curves (see \cite[Corollary~4.4]{Luo98:IntersectionNumbers}).
Feng Luo's work characterizes functions $f \colon \mathcal{S}(S) \to \mathbb{R}$ on simple closed curves that come from intersection numbers in terms of an infinite family of polynomial equations involving the $f$-value of triples of simple curves in the max-plus semi-ring structure on the real numbers.
A consequence of this characterization is that $h(n) \coloneqq f([a^n b])$ is a convex function.
In the proof of Corollary~\ref{cor:additivity-limit}, we generalize this convexity to possibly non-simple arbitrary curves $a,b$. See Proposition~\ref{cor:additivity-limit} to see how our conditions on a curve function $f$ imply Feng Luo's conditions for the restriction of $f$ to simple curves.
\end{remark}

 Here we show that all the necessary hypotheses are satisfied to apply Theorem~\ref{thm:caratheodory}.

\begin{proposition}\label{prop:join-split-alg}
Let $a,b,x$ be arbitrary elements of $\pi_1(S)$.
  If $a$ and $b$ cross $x$ to the right, then the limit
  \[
    \lim_{n,m\to \infty} f([a x^n b x^m]) - f([ax^n]) - f([bx^m])
  \]
  exists and is equal to~$0$.
\end{proposition}
\begin{proof}
We separate the proof in two cases, depending on the magnitude of $|m-n|$.
\begin{itemize}
    \item Case 1: $|n-m|$ is large.  By Proposition~\ref{prop:join-split-geom}, we have
  \[
    f([a x^n]) + f([b x^{m}]) \ge f([a x^{n+1} b x^{m-1}])
      \ge f([a x^{n+2}]) + f([b x^{m-2}]).
  \]
  for all $m,n$ so that $m-n -s >r > s$, or $m-n+s<r<-s$, i.e.
  For the first one, it suffices to have $m-n-s > s+1$, i.e., $m-n>2s+1$
  For the second, $m-n+s<-s-1$, i.e. $m-n<-2s-1$.
  In any case, $|m-n|>2s+1$
  Thus, by Corollary~\ref{cor:additivity-limit}, the claim follows in this case.
  \item Case 2: $|m-n|$ is small.  Let $p=m-n$ and suppose that $|p| \leq 2s+1$. By Lemma~\ref{lem:rhs_prod}  $ax^nb$ crosses $x$ to the right. Therefore, we have that for every $k$
  \[
  [ax^nbx^{n+p}] [x^k] \reducesto [ax^nbx^{n+p+k}].
  \]
  On the other hand, $bx^{n+p+k}a$ crosses $x$ to the right, also, and since
  \[
  [ax^{n+k}bx^{n+p+k}]=[x^{n+k}bx^{n+p+k}a]=[bx^{n+p+k}ax^{n+k}]
  \]
  we have
  \[
  [bx^{n+p}ax^{n}] [x^k] \reducesto [bx^{n+p}ax^{n+k}].
  \]  
  Therefore, by applying $f$ and using smoothing and additivity, we have that
  for all $k, n$
  \begin{align*}
  [ax^{n+p} bx^{n-k}] 2[x^k] & \reducesto  [ax^{n+p} bx^{n}] [x^k]  \\
  &\reducesto  [ax^{n+p} bx^{n+k}]
  \end{align*}
  Subtracting $[x^{p+2n+k}]$, applying $f$ using smoothing and additivity, and taking the limit when $n,k$ go to infinity we see, by Case 1, that
  \[
  \lim_{n \to \infty} f([ax^{n+p}bx^n])-f([x^n])=0
  \]  
  as we wanted to show.\qedhere
\end{itemize}
 \end{proof}

%%% Local Variables:
%%% mode: latex
%%% TeX-master: "Intersections"
%%% End:
